%!TEX root = ../main.tex
\chapter{Postcards from Nowhere}

The first postcard arrived with no return address and a postmark Elin
couldn't place. It showed a lighthouse she didn't recognize, and on
the back, in handwriting that was almost but not quite her brother's:
\emph{Made it further than I thought. More soon.}

She turned it over twice, looking for something she'd missed, then set
it on the kitchen table and forgot about it for three days, the way
you forget about anything that doesn't fit into the shape of your
week.

The second postcard came eleven days later. A different lighthouse,
or maybe the same one from a different angle. \emph{Weather's been
kind. Tell Mom I'm not dead, just slow.}

By the fourth postcard she had started keeping them in a shoebox, and
by the seventh she had stopped believing they were really from her
brother, who had, as far as anyone could confirm, not left the country
in four years and did not, to her knowledge, know anyone who owned a
lighthouse. But the handwriting was close enough, and the voice — the
particular dry understatement of it — was close enough, that she kept
reading them anyway, the way you keep answering a phone that might, on
some call, finally be who you're hoping for.

The last postcard, the one that came after the funeral, said only:
\emph{Almost there.}

She never found out where "there" was. She kept the shoebox anyway.
